\section{Chapitre~\ref{sec:debugging}. Déboguer la machine}

Ce qu'entend une électrode, la faible dimension de l'activité, le fonctionnement des interfaces à grilles rigides, l'apprentissage conjoint du cerveau et du décodeur, et les points visuels produits par stimulation sont mesurés dans des expériences publiées avec comité de lecture; les estimations des flux de données sont l'arithmétique du chapitre; sur l'appareil de Neuralink implanté chez l'humain, pour autant que nous ayons pu le vérifier, les données ont surtout été publiées par la société elle-même.

{\footnotesize
\setlength{\tabcolsep}{3pt}\renewcommand{\arraystretch}{1.15}
\begin{longtable}{>{\raggedright}p{0.5cm}>{\raggedright}p{4.0cm}>{\raggedright}p{5.6cm}>{\raggedright}p{2.3cm}>{\raggedright\arraybackslash}p{1.7cm}}
\toprule
N\textsuperscript{o} & Affirmation & Expérience et ce qui est mesuré & Source & Statut \\
\midrule\endfirsthead
\toprule
N\textsuperscript{o} & Affirmation & Expérience et ce qui est mesuré & Source & Statut \\
\midrule\endhead
1 & L'électrode entend les impulsions des neurones situés dans un rayon d'une centaine de micromètres, en général d'un seul à quelques-uns; on les distingue par leur forme & Rat, région CA1, enregistrement intracellulaire et extracellulaire simultané: la forme de l'impulsion extracellulaire reproduit la largeur et l'amplitude de l'impulsion intracellulaire. Estimation: une tétrode pourrait distinguer $60$--$100$ neurones, alors qu'en pratique on en isole environ six. & \cite{Henze2000extra} & mesuré \\
2 & Le cerveau compte $8.6\cdot10^{10}$ neurones; la sonde en entend environ un cent-millionième & Comptage des noyaux dans une suspension de tissu dissous: $86\cdot10^9$ neurones dans le cerveau humain. La fraction est l'arithmétique du chapitre. & \cite{Azevedo2009} & mesuré \\
3 & L'activité de centaines de neurones du cortex moteur tient dans une dizaine ou deux de variables communes & Revue d'enregistrements dans le cortex moteur du singe: l'activité de la population est décrite par quelques variables latentes communes à de nombreux neurones. & \cite{Gallego2017} & mesuré \\
4 & Le bruit des voisins est corrélé, et une centaine de neurones porte presque autant qu'un millier (proposition~\ref{prop:pooling}) & Cortex visuel du singe: coefficient de corrélation du bruit entre neurones voisins d'environ $0.1$, et le gain de la moyenne sature vers une centaine de neurones. Théorie des corrélations qui limitent l'information. & \cite{Zohary1994, MorenoBote2014} & mesuré \\
5 & Flux brut de $2\cdot10^8$~bit/s contre $8\cdot10^4$~bit/s en impulsions~\eqref{eq:probedata} & Arithmétique du chapitre à partir de la capacité du flux d'impulsions~\eqref{eq:spikeentropy}. Les paramètres de numérisation sont réels: la puce de Neuralink échantillonne à $19.3$~kHz sur $10$~bits par canal. & démonstration dans le chapitre; \cite{Rieke1997, Musk2019} & théorie \\
6 & Premier système Neuralink: les puces amplifient et numérisent, le flux brut part par câble, un programme externe détecte les impulsions; détection sur puce, boîtier fermé, radio et alimentation sans fil: pour l'appareil destiné à l'humain, selon la société & Article de la société: tout le flux brut passe par un câble USB-C, et un programme détecte les impulsions en temps réel hors de la puce; compression embarquée, alimentation et transmission sans fil sont présentées comme un projet pour les appareils cliniques. Pour l'appareil implanté chez l'humain, seulement des communications de la société. Que la détection des impulsions sur puce soit nécessaire est une conclusion du chapitre tirée de~\eqref{eq:probedata}. & \cite{Musk2019}; rapports de Neuralink, pas d'article à comité de lecture & mesuré (article de la société); chez l'humain: rapport de la société \\
7 & Jusqu'à $3072$ électrodes sur $96$ fils de $32$; un robot insère les fils en évitant les vaisseaux & Article de la société: $3072$ électrodes sur $96$ fils de $32$; le robot insère $6$ fils ($192$ électrodes) par minute; chez le rat porteur d'une grille implantée à long terme, jusqu'à $70\,\%$ des électrodes entendent des impulsions. & \cite{Musk2019} & mesuré (article de la société) \\
8 & Chez l'humain, environ mille canaux sur $64$ fils; une partie des fils s'est déplacée; curseur de l'ordre de $10$~bit/s & Communications de la société seulement. Une recherche dans PubMed n'a trouvé aucun article à comité de lecture présentant des résultats chez l'humain, ce qui concorde avec la réserve faite dans le texte du chapitre. & rapports de Neuralink; pas d'article à comité de lecture & mesuré (rapport de la société) \\
9 & Une interface sur une grille d'une centaine d'électrodes pilote un curseur & Personne paralysée, $96$ électrodes dans le cortex moteur primaire: l'intention de bouger le bras modifie les impulsions trois ans après la lésion; le décodeur fournit un \og curseur neuronal\fg{} (courrier, télévision), la commande d'une prothèse de main et d'un bras robotisé. & \cite{Hochberg2006} & mesuré \\
10 & Écriture \og par la main mentale\fg{}: environ $90$ caractères par minute, moins de $8$~bit/s & Personne à la main paralysée, tentatives d'écriture manuscrite, décodeur récurrent: $90$ caractères par minute avec une précision de $94.1\,\%$ en temps réel, plus de $99\,\%$ après correction automatique. L'estimation en bits est l'arithmétique du chapitre. & \cite{Willett2021} & mesuré \\
11 & Les tentatives de parole sont converties en texte à environ $60$ mots par minute & Personne ayant perdu une parole intelligible, grilles dans le cortex: $62$ mots par minute; $9.1\,\%$ d'erreurs sur les mots avec un vocabulaire de $50$ mots, $23.8\,\%$ avec un vocabulaire de $125\,000$. & \cite{Willett2023} & mesuré \\
12 & L'interface n'extrait qu'une petite part de la capacité: quelques dizaines de bits par seconde pour un neurone, quelques milliers pour une centaine (proposition~\ref{prop:probe}) & La borne supérieure~\eqref{eq:spikeentropy} est théorique; la comparaison avec les lignes~8 et~10 est une conclusion du chapitre. Que le goulot d'étranglement soit la faible dimension et la méconnaissance du code, et non le fil, n'a pas été vérifié par une expérience directe. & \cite{Rieke1997} & théorie \\
13 & Le cerveau et le décodeur apprennent l'un de l'autre & Singes pilotant un curseur; le décodeur s'adapte en boucle fermée pendant que les neurones réapprennent: la précision augmente, l'habileté se conserve et souffre moins des mouvements du bras propre. Le conseil de séparer les vitesses d'apprentissage est une règle d'ingénieur, sans expérience propre dans le chapitre. & \cite{Orsborn2014coad} & mesuré \\
14 & Le courant envoyé par une électrode excite les neurones et les fils environnants sans distinction de type & Cortex de souris et de chat, imagerie calcique à deux photons: même un faible courant excite des neurones épars, dispersés jusqu'à des millimètres, par excitation directe des axones dans un volume de quelques dizaines de micromètres de diamètre. L'excitation est donc non sélective, mais pas continue. & \cite{Histed2009stim} & mesuré \\
15 & La stimulation du cortex visuel allume des points; jusqu'à $16$ points à la fois permettent de reconnaître certaines lettres et les contours des objets & Personne aveugle, $96$ électrodes dans le cortex visuel pendant $6$ mois: seuil d'une électrode $66.8\pm36.5$~µA; la stimulation simultanée de groupes (jusqu'à $16$ électrodes, limite du stimulateur) abaisse le seuil et produit des motifs distincts; certaines lettres et les contours des objets sont reconnus. & \cite{Fernandez2021} & mesuré \\
16 & Le second axe de Neuralink est un appareil qui envoie des signaux dans le cortex visuel & Seulement des communications de la société sur son développement; aucune donnée à comité de lecture sur son fonctionnement n'a été trouvée. & rapports de Neuralink & hypothèse \\
17 & On peut mesurer les concentrations des neurotransmetteurs de diffusion par un capteur chimique dans le tissu & Tranches de cerveau de rat, voltampérométrie cyclique rapide sur fibre de carbone: libération et recapture de la sérotonine mesurées; la substance sort de la synapse. & \cite{Bunin1998} & mesuré \\
\bottomrule
\end{longtable}}
